\chapter{联想检索与算法研究}

\section{稳定问题合同与实验实现分离}

稳定抽象对应语义问题类别，例如 weak-cue discovery、bounded associative expansion 和 topology candidate generation；EPA、Residual、Wave、PPR 等具体算法只属于可替换实现/实验层。任何一个命名算法都不是 Nous Wave 的系统身份。比较替代方案时使用相同 Authority 输入、图构建规则、种子、硬过滤、计算预算、候选范围和诊断条件，先验证算法性质，再比较真实语料上的检索质量和总成本。\cite{nous-research-package}

\footnotesize
\begin{ReportWideTable}{>{\raggedright\arraybackslash}p{3.4cm}>{\raggedright\arraybackslash}p{3.7cm}>{\raggedright\arraybackslash}p{3.6cm}>{\raggedright\arraybackslash}p{3.2cm}}
\captionof{table}{联想检索机制及当前研究状态}\label{tab:algorithm-comparison}
\toprule
\textbf{机制} & \textbf{解决的问题} & \textbf{当前状态} & \textbf{下一步验证}\\
\midrule
Tag & 稳定语义线索身份。 & Authority primitive；同名字符串不等于同一 TagId。 & 身份解析、别名、aboutness 与 query semantics。\\
AssociativeTopology & 从类型化证据重建可查询图。 & Serving Projection；边导通率和中心性不是 Authority。 & 构建规则、fan/hub 抑制、重建一致性。\\
无拓扑基线 & 验证 graph expansion 是否真的增加有效召回。 & 必需对照。 & Exact/Entity/Lexical/Dense 与真实语料评测。\\
类型化 PPR/扩散 & 可解释、有界的图检索候选基线。 & 实验候选。 & 与无拓扑和 state-faithful Wave 同预算比较。\\
Wave & 有状态的竞争式联想扩展。 & 已实现研究分支；当前 state merge 存在预算近似，\textbf{未验证为默认方案}。 & 保留充分状态后与 PPR/无拓扑比较。\\
Residual Pyramid & 从当前 query residual 发现独立/弱语义方向。 & 数学核心局部差分较健康；作为实验 candidate generator。 & 测量端到端 recall/precision/cost。\\
EPA & 从 Tag vector 集构造全局语义基。 & 当前 weighted PCA 处理存在已知问题；实验性，默认不依赖。 & 修正 $\sqrt{weight}$ 协方差语义后再评估 utility。\\
固定计划融合 & 稳定融合多条检索证据。 & 目标基线为 fixed-plan weighted RRF；当前代码尚未对齐。 & 排序不变性与端到端检索评测。\\
\bottomrule
\end{ReportWideTable}
\normalsize

\section{关联 Authority 与 Serving 导通率}

AssociationEvidence 保存有类型和来源的关系证据，不直接保存一个能随使用次数无限累加的“最终 graph weight”。TopologyBuilder 按版本化规则将有效证据投影为当前图的 edge conductance，并记录 implementation id、revision 和 config digest。

图构建至少考虑：

\begin{itemize}
\item relation/association type；
\item polarity、validity 和 source independence；
\item fixed/normalized outbound mass；
\item fan/hub suppression；
\item query-dependent edge eligibility；
\item counterevidence 与生命周期状态。
\end{itemize}

被 retrieved、ranked、projected 或 exposed 本身不增加 Authority support。只有明确有独立语义的 Use/Outcome 或新 Observation 才能成为 owner 可接受的 evidence。这样检索系统不会通过自己的历史输出制造越来越强的自我确认回路。\cite{anderson1983,nous-research-package}

\section{Wave 与 PPR}

当前 Wave 实现将传播状态按前一节点和当前节点合并；不同路径如果具有不同 path budget，合并后的 energy 可以使用最强状态保留的预算继续传播，从而借用或丢失另一部分质量的可达性。该问题已经由构造/随机性质测试暴露，因此“代码存在”不等于“算法已经验证”。\cite{nous-code-snapshot,nous-research-package}

下一轮比较至少包含：

\begin{enumerate}
\item 无 topology expansion；
\item 类型化、有界 PPR/diffusion；
\item 保留影响后续可达性的充分状态的 Wave；
\item 如有必要，再加入 non-backtracking/single-hop 等较简单结构基线。
\end{enumerate}

所有方案共享相同 hard filters、种子和预算；同时记录截断质量、hub 影响、path/state 数、延迟和最终检索效用。Contradiction/negative relation 不应未经语义定义直接作为普通非负 transition。

\section{Residual Pyramid}

Residual 当前实现按 residual query 检索候选 Tag direction，经 modified Gram--Schmidt 构造局部子空间，将当前 residual 投影并相减，再依据剩余能量停止。研究资料包记录的 1000 组局部 differential test 在相同 exact-cosine candidate provider 下没有发现路径、能量或最终 residual 差异。该结果只支持核心数学转写在该测试条件下的一致性，不证明真实 retrieval quality。\cite{nous-research-package,nous-code-snapshot}

因此 Residual 可以保留为实验 candidate generator/query expansion seam，但不进入 Memory Reference Profile 的 correctness critical path。

\section{EPA 加权问题}

当前 EPA 在代表向量压缩后使用 multiplicity 计算 weighted mean，却在构造中心化 SVD matrix 时没有按 $\sqrt{weight}$ 缩放 row，因此没有保持原样本密度对应的 weighted covariance。研究材料中的构造反例显示主方向可发生约 $21.2^\circ$ 偏差，并在随机实验中频繁出现显著几何差异。\cite{nous-research-package,nous-code-snapshot}

这证明当前实现与目标 weighted PCA sample-density semantics 不一致，但不能直接推导 EPA 会降低最终检索质量。正确顺序是先修数学语义，再与更简单的 dense/multi-query/Residual baseline 比较；在此之前 EPA 保持实验性、默认不依赖。

\section{固定计划排名与拓扑通道}

当前 ranking 根据本批候选出现了哪些 EvidenceFamily 动态改变归一化分母，并在融合结果上额外叠加 topology continuous bonuses。前者已经有无关候选导致原候选相对顺序变化的反例，后者使 QueryPlan 难以解释最终排序。目标实现改为：

\begin{center}
QueryPlan freeze lanes/weights $\rightarrow$ lane ranking $\rightarrow$ object-level collapse $\rightarrow$ weighted RRF $\rightarrow$ optional rerank/diversity
\end{center}

Topology 与其他通道一样输出 object-level rank，不作为无法从计划解释的隐藏加分项。\cite{nous-research-package,nous-code-snapshot}

\section{阶段实现边界}

Memory Reference Profile 的可靠演示只要求 Exact、Entity、Lexical、Dense 等基础通道形成稳定闭环；高级 topology、EPA、Residual、Wave/PPR 均可在同一 semantic seam 后继续研究。这样算法实验不会决定 Memory Authority、Evidence、Revision、Use、Purge 等基础语义是否正确。